\section*{Glossary}
\label{sec:glossary}
\addcontentsline{toc}{section}{Glossary}

Reference definitions for recurring terms introduced throughout the paper.

\begin{description}[style=nextline]
  \item[Authentication]
    Verification of an asserted identity, origin, or authority through evidence such as credentials and signatures. Merkle integrity supplies the complementary structural check against a commitment.

  \item[Authorization]
    A policy decision determining whether a principal may perform an operation. In current federation, terminal-local rules grant remote principals path-scoped \texttt{get}, \texttt{set!}, and \texttt{resolve} permissions; administrators remain local principals.

  \item[Bridge]
    A reciprocal signed-head-retention relationship through which each journal preserves a verified head for the other. Reciprocity provides identity and history reachability, while terminal policy independently grants application-data authority.

  \item[Digest]
    A cryptographic hash derived from a node or larger structure and used as its integrity identifier.

  \item[Historical cursor]
    The independently selected origin and per-hop indexes used only by \texttt{resolve}. A canonical public journal-relative path interleaves these indexes with directional aliases; the Interface separates that path into the current route used for authentication and endpoint context and the historical cursor used to select committed content.

  \item[History]
    Used in its ordinary broad sense for prior states and changes. Concrete chains, heads, and receipts represent history using ordinary Synchronic Web structures.

  \item[Integrity]
    The property that modification can be detected relative to a known commitment. Authentication and authorization supply the associated claims about identity, origin, and authority.

  \item[Journal]
    The software an entity runs together with all data persisted on it. A journal supplies the root from which that entity views the Synchronic Web and is analogous to a site or origin on the conventional Web.

  \item[Journal identity]
    A stable commitment formed from a domain-separated random nonce and bound by bridges, preapprovals, and invocation audiences. Active signing keys may evolve beneath this identity through predecessor-authorized transitions.

  \item[Key-index]
    The terminal-local history index whose committed bridge state authenticated a remote interface key. Authorization may constrain this index relative to the terminal's current history.

  \item[Merge]
    To combine compatible partial structures whose integrity relationships agree.

  \item[Node]
    A generic term whose meaning follows from context. It may denote a node in the integrity-verifiable data-structure DAG or an entity in the graph of journals and relationships.

  \item[Object]
    A durable structure that couples executable behavior with state. Its documented protocol determines how the state is read, transformed, and represented; its internal layout remains implementation-specific.

  \item[Origin-local pin]
    Retention of verified remote proof material in the originating journal under the full origin-relative historical path. It follows a federated \texttt{resolve} and mutates the origin's retention state.

  \item[Path]
    A sequence interpreted from a selected journal root to address a resource or intermediate structure. A canonical federated path interleaves optional historical indexes with directional aliases before the terminal namespace and local path. Like a relative URL, its meaning is completed by a starting root.

  \item[Prune]
    To remove retained material while preserving the integrity relationships needed to identify the larger structure.

  \item[Record]
    Any snippet of data stored on a journal or transmitted between journals. The term is intentionally fluid across object shapes, historical entries, and storage units.

  \item[Resource]
    A path-addressed value or object in a journal-relative view. Its low-level representation may be a leaf or a larger self-encoding structure.

  \item[Root]
    A distinguished reference from which persisted structure is reached and interpreted. Unless otherwise qualified, a journal root is also the starting point for paths in that journal's view.

  \item[Root commitment]
    A digest or integrity reference accepted as the fixed point against which a structure is verified. Additional evidence is required to establish who created, supplied, or endorsed it.

  \item[Slice]
    A proof-preserving partial structure retaining the material required for selected paths while replacing unrelated material with stubs.

  \item[State]
    Used in its ordinary broad sense for the condition or contents of a system, structure, object, or resource at a relevant point or context.

  \item[Step]
    A discrete journal transition that turns staged changes into a new historically resolvable state.

  \item[Stub]
    A representation that retains the digest of unavailable, undisclosed, or intentionally pruned structure while omitting its contents. A stub represents unknown structure; null represents empty structure.

  \item[Stump]
    The low-level persisted node form used to represent a stub. It stores a retained logical digest under a distinct local node handle.

  \item[Synchronic Web]
    The network, data structure, computational medium, and protocol family formed when independently rooted journals make resources and histories reachable through journal-relative paths and bridges.

  \item[Terminal journal]
    The journal that directly receives and authorizes a federated application invocation after the origin follows a working route to its committed endpoint and authentication object.

  \item[Working route]
    The ordered receiver-local directional aliases used from the origin's latest committed view to select the terminal authentication object, endpoint, and reverse path to the origin key. Its implicit indexes are current; \texttt{resolve} may select historical content independently.

  \item[Entanglement]
    Cross-history ordering created when one journal incorporates another's committed head into its own later state. The term follows Maniatis and Baker's secure-history protocol~\cite{Maniatis2002timeline}.
\end{description}
